% Modified for NDSS 2024 by MS on 2024/07/01

%% bare_conf.tex
%% V1.4b
%% 2015/08/26
%% by Michael Shell
%% See:
%% http://www.michaelshell.org/
%% for current contact information.
%%
%% This is a skeleton file demonstrating the use of IEEEtran.cls
%% (requires IEEEtran.cls version 1.8b or later) with an IEEE
%% conference paper.
%%
%% Support sites:
%% http://www.michaelshell.org/tex/ieeetran/
%% http://www.ctan.org/pkg/ieeetran
%% and
%% http://www.ieee.org/

%%*************************************************************************
%% Legal Notice:
%% This code is offered as-is without any warranty either expressed or
%% implied; without even the implied warranty of MERCHANTABILITY or
%% FITNESS FOR A PARTICULAR PURPOSE! 
%% User assumes all risk.
%% In no event shall the IEEE or any contributor to this code be liable for
%% any damages or losses, including, but not limited to, incidental,
%% consequential, or any other damages, resulting from the use or misuse
%% of any information contained here.
%%
%% All comments are the opinions of their respective authors and are not
%% necessarily endorsed by the IEEE.
%%
%% This work is distributed under the LaTeX Project Public License (LPPL)
%% ( http://www.latex-project.org/ ) version 1.3, and may be freely used,
%% distributed and modified. A copy of the LPPL, version 1.3, is included
%% in the base LaTeX documentation of all distributions of LaTeX released
%% 2003/12/01 or later.
%% Retain all contribution notices and credits.
%% ** Modified files should be clearly indicated as such, including  **
%% ** renaming them and changing author support contact information. **
%%*************************************************************************


% *** Authors should verify (and, if needed, correct) their LaTeX system  ***
% *** with the testflow diagnostic prior to trusting their LaTeX platform ***
% *** with production work. The IEEE's font choices and paper sizes can   ***
% *** trigger bugs that do not appear when using other class files.       ***                          ***
% The testflow support page is at:
% http://www.michaelshell.org/tex/testflow/



\documentclass[conference]{IEEEtran}

% --- 数学与字体包 ---
\usepackage{amsmath}
\usepackage{amsfonts,amssymb}
\usepackage{mathrsfs}

% --- 解决 IEEEtran 与 amsthm 的冲突 ---
% 必须在加载 amsthm 之前执行这两行，否则会报错
\let\proof\relax 
\let\endproof\relax
\usepackage{amsthm}

% --- 图形与表格包 ---
\usepackage{graphicx}
\usepackage{textcomp}
\usepackage{xcolor}
\usepackage{booktabs} % 表格三线表
\usepackage{multirow} % 表格多行合并
\usepackage{array}
\usepackage{csquotes}

% --- IEEE 推荐的子图包设置 ---
\usepackage[caption=false,font=normalsize,labelfont=sf,textfont=sf]{subfig}

% --- 算法包 ---
\usepackage{algorithm}
\usepackage{algorithmic}

% --- 其他工具包 ---
\usepackage{url}
\usepackage{verbatim}
\usepackage{cite}
\usepackage{balance} % 用于平衡最后一页的双栏
\usepackage{stfloats} % 用于跨栏长公式/图片底部对齐
\usepackage{hyperref} % 超链接，通常放在最后加载
\usepackage{tabularx, makecell}

\theoremstyle{definition} % 使得定义和例子的正文内容不使用斜体
\newtheorem{definition}{Definition}

\theoremstyle{plain} % 使得定理、引理的正文内容使用斜体（学术界标准）
\newtheorem{theorem}{Theorem}
\newtheorem{lemma}{Lemma}
\newtheorem{proposition}{Proposition}


% *** GRAPHICS RELATED PACKAGES ***
%
\ifCLASSINFOpdf
  % \usepackage[pdftex]{graphicx}
  % declare the path(s) where your graphic files are
  % \graphicspath{{../pdf/}{../jpeg/}}
  % and their extensions so you won't have to specify these with
  % every instance of \includegraphics
  % \DeclareGraphicsExtensions{.pdf,.jpeg,.png}
\else
  % or other class option (dvipsone, dvipdf, if not using dvips). graphicx
  % will default to the driver specified in the system graphics.cfg if no
  % driver is specified.
  % \usepackage[dvips]{graphicx}
  % declare the path(s) where your graphic files are
  % \graphicspath{{../eps/}}
  % and their extensions so you won't have to specify these with
  % every instance of \includegraphics
  % \DeclareGraphicsExtensions{.eps}
\fi
% graphicx was written by David Carlisle and Sebastian Rahtz. It is
% required if you want graphics, photos, etc. graphicx.sty is already
% installed on most LaTeX systems. The latest version and documentation
% can be obtained at: 
% http://www.ctan.org/pkg/graphicx
% Another good source of documentation is "Using Imported Graphics in
% LaTeX2e" by Keith Reckdahl which can be found at:
% http://www.ctan.org/pkg/epslatex
%
% latex, and pdflatex in dvi mode, support graphics in encapsulated
% postscript (.eps) format. pdflatex in pdf mode supports graphics
% in .pdf, .jpeg, .png and .mps (metapost) formats. Users should ensure
% that all non-photo figures use a vector format (.eps, .pdf, .mps) and
% not a bitmapped formats (.jpeg, .png). The IEEE frowns on bitmapped formats
% which can result in "jaggedy"/blurry rendering of lines and letters as
% well as large increases in file sizes.
%
% You can find documentation about the pdfTeX application at:
% http://www.tug.org/applications/pdftex


 \pagestyle{plain}



% correct bad hyphenation here
\hyphenation{op-tical net-works semi-conduc-tor}


\begin{document}
%
% paper title
% Titles are generally capitalized except for words such as a, an, and, as,
% at, but, by, for, in, nor, of, on, or, the, to and up, which are usually
% not capitalized unless they are the first or last word of the title.
% Linebreaks \\ can be used within to get better formatting as desired.
% Do not put math or special symbols in the title.
\title{Scalable and Adaptive Network Hardening via MaxSAT: A Probabilistic Framework Based on Attack-Defense Graphs}


% author names and affiliations
% use a multiple column layout for up to three different
% affiliations
\author{\IEEEauthorblockN{Anonymous authors}}
	
% conference papers do not typically use \thanks and this command
% is locked out in conference mode. If really needed, such as for
% the acknowledgment of grants, issue a \IEEEoverridecommandlockouts
% after \documentclass

% for over three affiliations, or if they all won't fit within the width
% of the page, use this alternative format:
% 
%\author{\IEEEauthorblockN{Michael Shell\IEEEauthorrefmark{1},
%Homer Simpson\IEEEauthorrefmark{2},
%James Kirk\IEEEauthorrefmark{3}, 
%Montgomery Scott\IEEEauthorrefmark{3} and
%Eldon Tyrell\IEEEauthorrefmark{4}}
%\IEEEauthorblockA{\IEEEauthorrefmark{1}School of Electrical and Computer Engineering\\
%Georgia Institute of Technology,
%Atlanta, Georgia 30332--0250\\ Email: see http://www.michaelshell.org/contact.html}
%\IEEEauthorblockA{\IEEEauthorrefmark{2}Twentieth Century Fox, Springfield, USA\\
%Email: homer@thesimpsons.com}
%\IEEEauthorblockA{\IEEEauthorrefmark{3}Starfleet Academy, San Francisco, California 96678-2391\\
%Telephone: (800) 555--1212, Fax: (888) 555--1212}
%\IEEEauthorblockA{\IEEEauthorrefmark{4}Tyrell Inc., 123 Replicant Street, Los Angeles, California 90210--4321}}




% use for special paper notices
%\IEEEspecialpapernotice{(Invited Paper)}


\IEEEoverridecommandlockouts
\makeatletter\def\@IEEEpubidpullup{6.5\baselineskip}\makeatother
\IEEEpubid{\parbox{\columnwidth}{
		Network and Distributed System Security (NDSS) Symposium 2026\\
		23 - 27 February 2026 , San Diego, CA, USA\\
		ISBN 979-8-9919276-8-0\\  
		https://dx.doi.org/10.14722/ndss.2026.[23$|$24]xxxx\\
		www.ndss-symposium.org
}
\hspace{\columnsep}\makebox[\columnwidth]{}}


% make the title area
\maketitle

% As a general rule, do not put math, special symbols or citations
% in the abstract
\begin{abstract}
Effective cybersecurity management for large-scale networks requires a multi-objective trade-off among minimizing attack risks, controlling defense implementation costs, and maximizing system operational benefits. While Attack Graphs (AGs) serve as a standard formalism for modeling multi-step attacks, determining optimal hardening strategies remains computationally intractable due to the explosion of attack path combinations. Existing methods exhibit critical limitations: as AGs scale, the exponential growth of the defense strategy space renders exact algorithms computationally prohibitive. Furthermore, integrating attacker capabilities into node risk assessments requires iterative Bayesian inference updates across varying defense configurations, which becomes computationally intractable in large-scale graphs. To overcome these limitations, this paper proposes a novel framework that systematically transforms the network hardening problem into a Maximum Satisfiability (MaxSAT) instance. We construct an Attack-Defense Graph and introduce an encoding mechanism that embeds attack causality, countermeasure costs, and exploitation likelihoods directly into Weighted Conjunctive Normal Form (WCNF) clauses by using local probability derivation. Instead of directly optimizing the global utility function, the MaxSAT solver optimizes a tractable surrogate objective induced by local conditional probabilities. \textcolor{red}{We show that the WCNF formulation preserves the causal structure of the graph, that MaxSAT exactly optimizes this surrogate objective, and how the surrogate relates to the belief propagation evaluation of global utility. Experiments further show competitive solution quality on small graphs, clearer quality and runtime advantages on medium graphs, scalability to large topologies, and adaptation to policy constraints and observed attack events.}
\end{abstract}


% no keywords




% For peer review papers, you can put extra information on the cover
% page as needed:
% \ifCLASSOPTIONpeerreview
% \begin{center} \bfseries EDICS Category: 3-BBND \end{center}
% \fi
%
% For peerreview papers, this IEEEtran command inserts a page break and
% creates the second title. It will be ignored for other modes.
\IEEEpeerreviewmaketitle



\section{Introduction} \label{Introduction}


With the rapid expansion of modern network infrastructures, ranging from large-scale enterprise networks to emerging Internet of Things (IoT) environments, the complexity of cybersecurity management has increased exponentially~\cite{cabana2021threat,xiao2025caspr,10382176}. The intricate interdependencies between hosts and services create subtle vulnerability paths that attackers can exploit to launch sophisticated multi-step attacks. However, analyzing these paths becomes increasingly intractable as network scale grows, leading to a combinatorial explosion of potential attack scenarios that renders manual inspection computationally infeasible. Furthermore, effective security management is not merely about blocking intrusion paths; it requires solving a complex multi-objective trade-off involving potential attack losses, system service benefits, and defense implementation costs~\cite{wang2025heimdall}. Faced with such intricate cost-benefit calculations, security administrators often struggle to identify optimal defense strategies that maximize overall utility. This challenge is further compounded by the evolving nature of threats, necessitating automated frameworks capable of not only static optimization but also dynamic response to real-time changes.

\begin{figure}[htbp]
    \centering
    \includegraphics[width=0.95\linewidth]{flowchart.pdf} % 请确保文件名与您上传的一致
    \caption{The overall workflow of the proposed MaxSAT-based defense hardening framework. The system transforms network configurations into an Attack-Defense Graph, encodes it into a Boolean formula, and solves for the optimal defense strategy.}
    \label{fig:framework}
\end{figure}

To systematically analyze these threats, Attack Graphs (AGs) have been widely adopted as a standard modeling formalism~\cite{11187396, noel2003efficient,noel2004correlating}. Formally, an AG is represented as a directed graph where nodes characterize critical system states, including host privileges, attacker exploitation rules, network connectivity, and specific system vulnerabilities~\cite{alhomidi2012attack}. The edges connecting these nodes signify causal dependencies, often embodying logical conjunctive (AND) or disjunctive (OR) relationships that dictate how an atomic attack step enables subsequent exploits. By capturing these complex interdependencies, AGs enable administrators to visualize potential intrusion scenarios. However, as network scale grows, the analysis of attack graphs faces the problem of combinatorial explosion, making the manual determination of optimal defense strategies computationally infeasible. 

Despite the extensive literature on attack graph-based security analysis, existing methodologies exhibit distinct limitations when applied to practical hardening scenarios~\cite{dewri2007optimal,mehta2006ranking}. A significant portion of current research is dedicated to critical vulnerability identification and critical path analysis. While these approaches excel at diagnosing topological weaknesses, they frequently abstract away the operational constraints of defense deployment, failing to translate abstract risk scores into actionable, compatible combinations of countermeasures. Conversely, frameworks that explicitly model defense selection face significant hurdles. The introduction of defensive countermeasures dynamically alters network reachability, necessitating a computationally prohibitive re-evaluation of probabilistic inference that severely impedes real-time adaptability. Moreover, the reliance on heuristic search algorithms fails to guarantee convergence to a global optimum, often yielding sub-optimal strategies insufficient for high-stakes security environments. Furthermore, even with efficient heuristic searches, evaluating attacker success probabilities requires re-executing probabilistic inference for every candidate defense. Although methods like Loopy Belief Propagation (Loopy-BP) operate in polynomial time, the exponential explosion of defense combinations across expanding attack graphs renders this repetitive computation intractable.

Among prior studies, the framework proposed by Poolsappasit et al.~\cite{poolsappasit2011dynamic} is particularly relevant to our work, as it also considers dynamic security risk management under probabilistic attack dependencies and jointly accounts for security risks, operational benefits, and defense costs. While our research shares a similar problem setting with this line of inquiry, it diverges fundamentally in methodology. Rather than combining Bayesian-graph-based probabilistic evaluation with heuristic search (e.g., Genetic Algorithms), we reformulate defense selection on Attack-Defense Graphs as a MaxSAT optimization problem. This reformulation yields a formal encoding of logical dependencies, deployment constraints, and a tractable surrogate objective derived from local probabilities. As shown in our evaluation, this shift in optimization paradigm delivers substantially better scalability while maintaining strong solution quality, especially on large network instances.


Guided by this formal reformulation, our proposed defense hardening framework, as illustrated in Fig.~\ref{fig:framework}, systematically integrates attack causality, exploitation likelihoods, and defensive countermeasures. We first introduce an Attack-Defense Graph (ADG), constructed by augmenting standard attack graphs with defense nodes that explicitly model the interaction between security countermeasures and network configurations. This formalism allows us to bridge the gap between risk assessment and constraint optimization by mapping the network hardening problem directly to a Maximum Satisfiability (MaxSAT) instance. Unlike traditional methods that separate probabilistic inference from logical solving, our approach systematically embeds the attackers' success probabilities directly into the WCNF clauses using local probability derivation. By leveraging state-of-the-art MaxSAT solvers, our framework efficiently computes high-quality defense portfolios by exactly optimizing the formulated tractable metric that balances security benefits against implementation costs. Furthermore, we provide a theoretical proof of the equivalence between our logic-based formulation and the risk optimization objective. The anonymized artifact package, including our implementation, WCNF encoder, benchmark generator, retained MulVAL input files, the complete generated attack graph, and scripts for reproducing the evaluation, is available at \url{https://anonymous.4open.science/r/SAT_pro-8FF0/README.md}.

The main contributions of this paper are summarized as follows:

{\color{red}
\begin{itemize}
    \item[i)] We formulate network hardening as a defender-oriented Weighted Partial MaxSAT problem over an Attack-Defense Graph. Hard clauses encode attack causality and shared defense mappings, while soft clauses represent operational benefit, deployment cost, and probabilistic local risk.

    \item[ii)] We develop a WCNF encoding that avoids repeated global probabilistic inference during portfolio search. Theoretical analysis proves that the encoding preserves causal dependencies and that MaxSAT exactly optimizes the encoded surrogate objective, and further characterizes its relation to the BP evaluation of global utility.
    
    \item[iii)] We evaluate the framework on a VM-based enterprise testbed and synthetic graphs. The results demonstrate its solution quality, computational efficiency, dynamic adaptability, and robustness under varying graph structures and model parameters.
\end{itemize}
}

The remainder of this paper is organized as follows. Section~\ref{Related Work} reviews related work. Section~\ref{sec:system_modeling} introduces our system modeling approach, detailing the formal generation of Attack-Defense Graphs and the hardening objective function. Section~\ref{sec:methodology} elaborates on the proposed methodology, describing the WCNF encoding and theoretical proofs.  Section~\ref{sec:experiments} presents a comprehensive evaluation of the proposed framework, encompassing effectiveness case studies, dynamic adaptability assessments, and extensive comparative and scalability analyses. Finally, Section~\ref{sec:conclusion} concludes the paper.

{\color{red}
\section{Related Work}\label{Related Work}

\subsection{Attack Graph Generation}

Phillips and Swiler \cite{phillips1998graph} established an early framework for modeling multistep intrusions, while Schneier \cite{schneier1999attack} introduced attack trees. Sheyner et al. \cite{sheyner2002automated} subsequently modeled attacks through explicit state enumeration, but the resulting state-space explosion limited scalability. Ammann et al. \cite{10.1145/586110.586140} addressed this limitation through the monotonicity assumption, enabling exploit-dependency representations that grow polynomially with network size. Jha et al. \cite{jha2002two} further developed formal analyses of attack dependencies.

Ou et al. \cite{ou2005mulval,ou2006scalable} developed MulVAL, which encodes network configurations, vulnerabilities, and exploitation rules in Datalog \cite{ceri1989you} to generate logical attack graphs automatically. Host-based abstractions \cite{ammann2005host,ingols2006practical} further reduce graph size by representing host reachability and privilege levels, although they omit fine-grained exploit dependencies. To support quantitative analysis, Liu and Man \cite{liu2005network} combined attack graphs with Bayesian networks, providing a probabilistic basis for subsequent risk assessment and network-hardening research.

\subsection{Network Hardening and Risk Assessment}

Jha et al. \cite{jha2002two} formulated network hardening as a Minimum Hitting Set problem and showed that finding a minimum exploit set that disrupts all attack paths is NP-complete. Subsequent studies therefore employed approximation and heuristic methods. Chao et al. \cite{chao2014heuristic} identified minimal sets of initial conditions, Alhomidi and Reed \cite{alhomidi2012attack} applied genetic algorithms to minimum-cut analysis, and Al Ghazo et al. \cite{al2023critical} used strongly connected components to improve critical-cut identification.

Probabilistic methods additionally account for uncertainty in exploit success. Poolsappasit et al. \cite{poolsappasit2011dynamic} combined Bayesian attack graphs, CVSS-derived likelihoods, and genetic algorithms to guide defense selection. Bansal et al. \cite{bansal2023quantitative} incorporated vulnerability complexity into cost and benefit analysis to reduce unnecessary patching. Exact probability inference has also been studied: Homer et al. \cite{homer2009sound} proposed a sound iterative method whose exponential complexity limits scalability, while Mu\~{n}oz-Gonz\'{a}lez et al. \cite{munoz2017exact} compared exact and approximate inference methods in terms of accuracy and efficiency.

Recent optimization-based studies address richer hardening objectives. Khouzani et al. \cite{KHOUZANI2019894} derive an exact min--max MILP that jointly considers risk and control costs, including controls affecting multiple vulnerabilities. Zenitani \cite{zenitani2022multi} uses monotonic metrics to approximate Pareto-efficient countermeasure sets, while Wan et al. \cite{wan2023iot} combine shortest-attack-trace analysis with exact and heuristic hardening algorithms. These methods rely respectively on MILP, monotonic search, and graph-specific algorithms.

\subsection{SAT-based Security Management}

Prior logic-based methods use Boolean formulations for different security decisions. Wang et al. \cite{wang2006minimum} expand a protection formula into DNF to derive minimum-cost hardening sets, with potentially exponential growth. Homer and Ou \cite{homer2009sat} construct a linear-size CNF and use MinCostSAT to balance reconfiguration and attack-damage costs, while UNSAT-core analysis identifies conflicts between security and usability policies. Huang et al. \cite{huang2011distilling} assign attacker-side costs to CNF literals and iteratively extract critical attack paths and surfaces. 

Wang et al. and Homer and Ou retain deterministic reachability semantics, while Huang et al. use exploit likelihoods to rank paths from the attacker perspective rather than to define a residual risk objective. The reviewed formulations do not jointly model shared defense actions, preserved operational benefit, deployment cost, and a risk surrogate derived from local probabilities. Our Weighted Partial MaxSAT (WPM) formulation introduces shared defense variables and hard clauses that encode attack causality, policy constraints, and mappings from each defense to multiple configurations. Weighted soft clauses represent preserved operational benefit, deployment cost, and risk penalties derived from local probabilities. A MaxSAT solver exactly optimizes this tractable surrogate within a single WPM instance, rather than repeatedly evaluating global probabilistic utility for candidate defense portfolios. Table~\ref{tab} compares our formulation with representative SAT-based attack graph methods.
\begin{table*}[t]
\centering
\caption{Comparison with prior logic- and SAT-based attack-graph methods.}
\label{tab}
\footnotesize
\setlength{\tabcolsep}{3pt}
\renewcommand{\arraystretch}{1.12}

\begin{tabular}{
@{}
>{\raggedright\arraybackslash}m{3.25cm}
>{\raggedright\arraybackslash}m{2.90cm}
>{\raggedright\arraybackslash}m{3.40cm}
>{\raggedright\arraybackslash}m{3.40cm}
>{\raggedright\arraybackslash}m{3.20cm}
@{}
}
\toprule
\textbf{Work / Encoding} &
\textbf{Output} &
\textbf{Risk Treatment} &
\textbf{Defense Abstraction} &
\textbf{Optimized Terms} \\
\midrule

Wang et al. (DNF) \cite{wang2006minimum} &
\makecell[l]{Minimum-cost\\hardening set} &
Deterministic reachability &
Initial conditions &
Hardening cost \\

\addlinespace[2pt]

Homer--Ou (CNF) \cite{homer2009sat} &
Secure configuration &
Deterministic reachability &
Configuration facts &
\makecell[l]{Configuration, usability,\\and damage costs} \\

\addlinespace[2pt]

Huang et al. (CNF) \cite{huang2011distilling} &
Critical attack surface &
\makecell[l]{Likelihood-based path\\ranking} &
No defense selection &
Attacker effort \\

\midrule

\textbf{Our work (WCNF)} &
Defense portfolio &
\makecell[l]{Local probability-weighted\\risk} &
\makecell[l]{Shared defense variables\\(one-to-many)} &
\makecell[l]{Benefit, deployment cost,\\and risk} \\

\bottomrule
\end{tabular}
\end{table*}

}
\section{System Modeling}
\label{sec:system_modeling}

In this section, we formally define the network attack graph model to explicitly characterize the causal relationships between vulnerability exploitations and network states. Beyond the fundamental definition, we extend the graph structure to incorporate defense measures, enabling a unified representation of attack paths and mitigation strategies. Furthermore, we propose a comprehensive security metric formulation that quantifies the trade-offs among system operational benefits, defense implementation costs, and potential security risks. Consequently, the hardening problem is modeled as an optimization problem aimed at maximizing the overall system utility. For convenience, the key notations and symbols used throughout this paper are summarized in Table~\ref{tab:notations}.

\subsection{Formal Modeling of Attack-Defense Dependencies}
\label{subsec:adg_modeling}

\begin{table}[htbp]
    \centering
    % 1. 关键修改：将行间距设置为默认值的 1.3 倍，显著提升可读性
    \renewcommand{\arraystretch}{1.3} 
    \caption{Summary of Important Notations}
    \label{tab:notations}
    
    % 2. 这里的宽度建议：符号列给 0.12-0.15，描述列给 0.8-0.85
    % 确保总和不超过 1 (考虑 tabcolsep 边距，通常预留一点缓冲)
    \begin{tabular}{p{0.12\columnwidth} p{0.75\columnwidth}}
        \toprule
        \textbf{Symbol} & \textbf{Description} \\
        \midrule
        
        % Group 1: Graph Structures & Sets
        $\mathcal{G}_{base}$ & The logical attack graph structure \\
        $\mathcal{G}$        & The attack-defense graph \\
        $N_p$                & The set of primitive fact nodes $c$ \\
        $N_d$                & The set of derived fact nodes $p$ \\
        $N_r$                & The set of exploit nodes $e$ \\
        $D$                  & The set of available defense countermeasures $d$ \\
        
        \addlinespace % 添加逻辑分组间距
        
        % Group 2: Vectors
        $\mathbf{x}$         & The binary defense decision vector \\
        $\mathbf{v}$ & The binary vector denoting the status of child exploits \\
        $\mathbf{s}$         & The global state vector encompassing the truth values of all node types ($c, p, e, d$) \\
        
        \addlinespace % 添加逻辑分组间距
        
        % Group 3: Parameters (Benefits & Costs)
        $B_c$                & The benefit of retaining configuration $c$ \\
        $C_d$                & The deployment cost of defense countermeasure $d$ \\
        $I_p$                & The impact value of compromising privilege $p$ \\
        
        \addlinespace % 添加逻辑分组间距
        
        % Group 4: Probabilities & Penalties
        
        {\color{red}$\rho_e$} & {\color{red}The conditional success probability of exploit $e$ given that all its preconditions are satisfied} \\
        $\Pr(p|\mathbf{v})$ & The probability of compromising privilege $p$ given the activation state vector $\mathbf{v}$ of child exploits \\
        $\mathcal{W}_{p,\mathbf{v}}$ & The penalty cost associated with privilege $p$ given the activation state vector $\mathbf{v}$ of child exploits \\
        
        \bottomrule
    \end{tabular}
\end{table}

To facilitate the rigorous optimization of network hardening strategies, we establish a formal graph-theoretic framework in this section. We begin by leveraging the standard Logical Attack Graph (LAG) to characterize the adversary's lateral movement capabilities and causal dependencies. Subsequently, we augment this foundation by explicitly integrating defense measures into the graph structure, resulting in the proposed Attack-Defense Graph. This unified representation serves as the basis for translating the hardening problem into a mathematical constraint satisfaction model. 

\subsubsection{Logical Attack Graph (Base Model)}
To systematically capture the causal dependencies between network configurations and potential adversary privileges, we leverage the Logical Attack Graph formalism proposed by Homer and Ou et al. \cite{ou2005mulval, ou2006scalable}. This model is generated by the logic-based network security analyzer, MulVAL, which utilizes Datalog derivation rules to reason about multi-host, multi-stage attacks.

\begin{definition}[Logical Attack Graph]
A Logical Attack Graph is defined as a six tuple $\mathcal{G}_{base} = (N_r, N_p, N_d, E, \mathcal{L}, G)$, where:
\begin{itemize}
    \item $N_r$, $N_p$, and $N_d$ are three disjoint sets of nodes representing \textit{derivation nodes}, \textit{primitive fact nodes}, and \textit{derived fact nodes}, respectively.
    \item $E \subseteq (N_r \times (N_p \cup N_d)) \cup (N_d \times N_r)$ denotes the set of directed edges representing causal dependencies.
    \item $\mathcal{L}$ is a mapping from a node to its semantic label (e.g., vulnerability description or host privilege).
    \item $G \subseteq N_d$ denotes the specific attacker's goal states.
\end{itemize}
\end{definition}

In this formalism, primitive fact nodes ($N_p$) encapsulate the initial network state (e.g., existence of vulnerabilities, active services, and connectivity), while derived fact nodes ($N_d$) represent the privileges acquired by the attacker through lateral movement. Derivation nodes ($N_r$) signify the logic rules or exploit actions that enable transitions between states.

To facilitate the subsequent SAT-based formulation and maintain consistency with the notation used in \cite{homer2009sat}, we map these formal sets to simpler variable types—exploits ($e$), configurations ($c$), and privileges ($p$)—as follows:

\begin{itemize}
    \item \textbf{Exploit Nodes ($e$):} Correspond to the derivation nodes ($N_r$). An exploit node $e \in N_r$ represents a specific attack action. It follows an AND-logic, requiring all preconditions (child nodes) to be satisfied. {\color{red}Each exploit node is associated with a parameter ($\rho_e\in(0,1]$), denoting the conditional success probability assumed by the model once all modeled prerequisites are satisfied.}
    \item \textbf{Configuration Nodes ($c$):} Correspond to the primitive fact nodes ($N_p$). A configuration node $c \in N_p$ represents atomic network facts or vulnerabilities. These serve as the leaf nodes and fundamental preconditions in the graph.
    \item \textbf{Privilege Nodes ($p$):} \textcolor{red}{Correspond to the derived fact nodes ($N_d$).} A privilege node \textcolor{red}{$p \in N_d$} represents an intermediate or final capability gained by the attacker. It follows an OR-logic, implying the privilege is acquired if any one of the outgoing exploits succeeds.
\end{itemize}

Fig. \ref{fig:original attack graph} illustrates an example of a standard Logical Attack Graph generated by MulVAL. In this graphical representation, the diamond-shaped nodes correspond to privileges ($p$), the rectangular nodes denote configurations ($c$), and the elliptical nodes represent exploits ($e$).

\begin{figure}[htbp]
    \centering
    % height=5cm：强制高度为5厘米
    % width=\linewidth：尝试占满单栏宽
    % keepaspectratio：保证图片不被拉伸变形（此时宽度会自动变窄以适应高度）
    \includegraphics[height=4cm, width=\linewidth, keepaspectratio]{original attack graph.pdf}
    \caption{An example of a Logical Attack Graph generated by MulVAL, illustrating the dependency between exploits and network configurations.}
    \label{fig:original attack graph}
\end{figure}

\subsubsection{Attack-Defense Graph Model}
\label{subsec:adg_model}

While the standard Logical Attack Graph effectively models how an attack propagates, it lacks an explicit representation of how to interrupt these paths within the graph structure itself. Merely identifying critical paths is insufficient for hardening optimization, necessitating the integration of feasible mitigation actions (e.g.,\ patching a vulnerability or modifying a firewall rule) and their logical relationships to attack preconditions. To address this limitation and enable SAT-based solvers to find optimal defense strategy, we extend the LAG to an Attack-Defense Graph.

{\color{red}
\begin{definition}[Attack-Defense Graph]
The Attack-Defense Graph is defined as a directed graph $\mathcal{G}=(\mathcal{V}_{ADG},\mathcal{E}_{ADG},\rho)$, constructed by augmenting the base LAG components $(\mathcal{V}_{base},\mathcal{E}_{base})$ with a set of defense nodes and defense edges.
\begin{equation}
    \mathcal{V}_{ADG}=\mathcal{V}_{base}\cup D,
\end{equation}
\begin{equation}
    \mathcal{E}_{ADG}=\mathcal{E}_{base}\cup\mathcal{E}_{def}.
\end{equation}
Here, $D$ represents the set of defense measures, and $\mathcal{E}_{def}\subseteq N_p\times D$ represents the set of directed edges connecting configuration nodes to their corresponding defense nodes. The mapping $\rho:N_r\rightarrow(0,1]$ assigns each exploit node $e$ a conditional success probability $\rho_e$ given that all its preconditions are satisfied.
\end{definition}
}

This extension preserves the fundamental causal logic of the original attack graph while incorporating remediation capabilities. In the base LAG, edges represent the dependencies required for an attack to succeed (i.e., ``if configuration $C$ exists, exploit $E$ becomes possible"). Analogously, the newly added edges in $\mathcal{E}_{def}$ represent the dependencies required for defense (i.e., ``if the administrator intends to invalidate configuration $C$, defense countermeasure $D$ must be adopted").

As illustrated in Fig. \ref{fig:ADG}, the Defense Nodes ($D$), depicted as trapezoids, act as logical inhibitors. The interaction mechanism follows a strict Inhibition Logic:
\begin{equation}
    d \Rightarrow \neg c, \quad \forall (c, d) \in \mathcal{E}_{def},
\label{eq:dc}
\end{equation}
where $c \in N_p$ is a mutable configuration and $d \in D$ is the corresponding defense. This logical implication signifies that the active deployment of a defense countermeasure ($d=1$) effectively negates the targeted precondition ($c=0$), thereby severing any causal chains dependent on $c$. By explicitly embedding these actions within the graph's causal structure, we can systematically translate the ADG into Conjunctive Normal Form (CNF) clauses. This transformation is pivotal, as it allows us to model the hardening problem as a Maximum Satisfiability (MAXSAT) instance, enabling the rigorous optimization of defense strategies based on the defined logical constraints.

% 插入 Fig 2：修改后的 ADG 图
\begin{figure}[htbp]
    \centering
    \includegraphics[height=4.5cm, width=\linewidth, keepaspectratio]{ADG.pdf}
    \caption{The proposed Attack-Defense Graph structure. Defense nodes (trapezoids) are introduced to inhibit primitive fact nodes. The directed edges from primitive facts to defense nodes indicate the feasible mitigation measures available to invalidate a specific network configuration. For instance, the edge $c_1 \to d_1$ signifies that the primitive fact $c_1$ can be invalidated by implementing defense $d_1$.}
    \label{fig:ADG}
\end{figure}

\subsection{Hardening Optimization Formulation}
\label{subsec:optimization}

The ultimate goal of our modeling framework is to determine the optimal defense strategy that maximizes system utility. We define a binary decision vector $\mathbf{x} = \{x_d \mid d \in D\}$, where $x_d = 1$ indicates that the defense countermeasure $d$ is deployed, and $x_d = 0$ otherwise.

We formulate the hardening problem as a global utility maximization problem:
\begin{equation}
    \max_{\mathbf{x}}\, \mathcal{U}(\mathbf{x}) = \mathcal{B}(\mathbf{x}) - \mathcal{C}(\mathbf{x}) - \mathcal{R}(\mathbf{x}),
\label{eq:global_function}
\end{equation}
where the objective function components are defined as follows:

\subsubsection*{1) System Operational Benefits ($\mathcal{B}$)}
Network services and configurations generate value for the organization. We assume that a mutable primitive fact node $c \in N_p$ provides a benefit $B_c$ only if it remains active (i.e., not inhibited by any deployed defense).
\begin{equation}
    \mathcal{B}(\mathbf{x}) = \sum_{c \in N_p} \mathbb{I}_c(\mathbf{x}) \cdot B_c.
\end{equation}
Here, $\mathbb{I}_c(\mathbf{x})$ is a state indicator function derived from the inhibition logic defined in Eq. \eqref{eq:dc}. Specifically, $\mathbb{I}_c(\mathbf{x}) = 1$ if the configuration $c$ is preserved (i.e., $\forall (c, d)\in \mathcal{E}_{def}, x_d = 0$), and 0 otherwise.

\subsubsection*{2) Defense Implementation Costs ($\mathcal{C}$)}
Deploying defense measures incurs explicit resource consumption (e.g., operational downtime or labor costs). The total implementation cost is the sum of costs for all active defense nodes:
\begin{equation}
    \mathcal{C}(\mathbf{x}) = \sum_{d \in D} x_d \cdot C_d,
\end{equation}
where $C_d$ represents the cost associated with implementing defense countermeasure $d$.

\subsubsection*{3) Residual Security Risks ($\mathcal{R}$)}
The security risk quantifies the expected impact of successful attacks under the defense strategy $\mathbf{x}$. This is calculated by aggregating the risk exposure of all privilege nodes:
\begin{equation}
    \mathcal{R}(\mathbf{x}) = \sum_{\textcolor{red}{p \in N_d}} \Pr(p \mid \mathbf{x}) \cdot I_p,
\end{equation}
where $I_p$ denotes the impact (or severity) of the attacker gaining privilege $p$, and $\Pr(p \mid \mathbf{x})$ represents the probability that the attacker successfully compromises node $p$ given the defense strategy $\mathbf{x}$. By minimizing this term, the optimizer seeks to suppress the most probable and damaging attack paths.

\section{Methodology}
\label{sec:methodology}

This section presents a logic-based computational framework to solve the optimal hardening problem defined in Section \ref{sec:system_modeling}. By leveraging the structural properties of the Attack-Defense Graph, we rigorously transform the global utility maximization problem into a Weighted Partial MaxSAT (WPM) instance. This transformation bridges the gap between domain-specific security semantics and general-purpose logical reasoning, enabling the use of efficient SAT solvers to navigate the complex trade-off space between security risks, operational benefits, and defense costs.

\subsection{Transformation of ADG into Weighted Logic Formula}
\label{subsec:transformation}

Formally, a WPM instance $\Phi$ is defined as a multiset of Boolean clauses partitioned into two disjoint sets: \textit{hard clauses} ($\Phi_{hard}$) and \textit{soft clauses} ($\Phi_{soft}$).
\begin{equation}
    \Phi = \Phi_{hard} \cup \Phi_{soft}.
\end{equation}
The hard clauses encode the inviolable topological constraints of the network, assigned with an infinite weight. The soft clauses encode the competing optimization objectives, assigned with finite weights. The goal is to find a truth assignment that satisfies all clauses in $\Phi_{hard}$ and maximizes the sum of weights of satisfied clauses in $\Phi_{soft}$.

\subsubsection{Topological Constraints (Hard Clauses)}
The construction of $\Phi_{hard}$ establishes a logical isomorphism between the graph connectivity and Boolean satisfiability constraints. We associate a Boolean variable with each node $v \in \mathcal{V}_{ADG}$. Logical truth values propagate bottom-up, reflecting the dependency edges from preconditions to goals.

\textbf{Exploit Nodes (Conjunctive Dependencies):}
An exploit node $e \in N_r$ follows an AND-logic structure. It is active if and only if all preconditions (child nodes) are satisfied. We formulate this strict dependency as:
\begin{equation}
    \label{eq:hard_exploit}
    C(e) = \bigvee_{v \in child(e)} \neg v \vee e,
\end{equation}
where $child(e)$ denotes the set of preconditions for exploit $e$. This clause forces $e$ to be true when all preconditions are true.

\textbf{Privilege Nodes (Disjunctive Dependencies):}
A privilege node \textcolor{red}{$p \in N_d$} follows an OR-logic structure. It is compromised if at least one outgoing exploit (child node) succeeds. The encoding enforces logical equivalence:
\begin{equation}
    \label{eq:hard_privilege}
    C(p) = \left( \neg p \vee \bigvee_{e \in child(p)} e \right) \wedge \bigwedge_{e \in child(p)} (\neg e \vee p).
\end{equation}

\textbf{Configuration Nodes (Inhibition Logic):}
    Configuration nodes $c \in N_p$ are leaf nodes representing network facts, conditioned on the absence of defense measures. Let $child(c) \subseteq D$ be the set of inhibiting defenses. The validity of $c$ is encoded as:
\begin{equation}
    \label{eq:hard_defense}
    C(c) = \bigwedge_{d \in child(c)} (\neg c \vee \neg d) \wedge \left( c \vee \bigvee_{d \in child(c)} d \right).
\end{equation}

The complete set of hard clauses is the conjunction of the above constraints:
\begin{equation}
    \Phi_{hard} = \bigwedge_{e} C(e) \wedge \bigwedge_{p} C(p) \wedge \bigwedge_{c} C(c).
\end{equation}

\subsubsection{Optimization Objectives (Soft Clauses)}
To rigorously map the utility function $\mathcal{U}(\mathbf{x})$ into the WPM format, we translate economic and security objectives into Boolean preferences.

\textbf{Benefit Preservation:}
Network services generate value only when accessible. For each configuration node $c$ with business value $B_c$, we introduce a soft clause $(c)$ with weight $w = B_c$. This incentivizes the solver to keep services active (i.e., not inhibited by defense).

\textbf{Cost Avoidance:}
Defense measures consume resources. For each defense node $d$ with implementation cost $C_d$, we introduce a soft clause $(\neg d)$ with weight $w = C_d$. This clause penalizes the activation of defense measures.

\textbf{Probabilistic Risk Mitigation:}
For a privilege node $p$ with immediate exploit children $\operatorname{child}(p)=\{e_1,\ldots,e_k\}$, let $\mathbf{v}\in\{0,1\}^{k}$ denote their Boolean activation states, where $v_i=1$ means that exploit $e_i$ is active. Assuming conditionally independent exploit outcomes, the compromise probability is
\begin{equation}
\Pr(p\mid\mathbf{v})
=
1-\prod_{i=1}^{k}(1-\rho_{e_i})^{v_i},
\label{eq:prob}
\end{equation}
where $\rho_{e_i}$ is the conditional success probability of $e_i$ once its prerequisites are satisfied. For each non-zero state $\mathbf{v}$, we construct
\begin{equation}
C_{\mathbf{v}}
=
\bigvee_{i=1}^{k}L(e_i,v_i),
\qquad
L(e_i,v_i)
=
\begin{cases}
\neg e_i, & v_i=1,\\
e_i, & v_i=0.
\end{cases}
\label{eq:risk_clause}
\end{equation}
The clause weight is $\mathcal{W}_{p,\mathbf{v}}=\Pr(p\mid\mathbf{v})I_p$. Since clauses are generated only for non-zero state vectors, exactly one matching blocking clause is falsified for each modeled non-zero exploit-state assignment, and the corresponding risk penalty is incurred exactly once.

Finally, the comprehensive set of soft clauses is synthesized as follows:
\begin{equation}
\begin{aligned}
    \Phi_{soft} = &\underbrace{\bigcup_{c} \{(c, B_c)\}}_{\text{Benefit}} \cup \underbrace{\bigcup_{d} \{(\neg d, C_d)\}}_{\text{Cost}} \\
    &\cup \underbrace{\bigcup_{p} \bigcup_{\mathbf{v} \neq \mathbf{0}} \{(C_{\mathbf{v}}, \mathcal{W}_{p,\mathbf{v}})\}}_{\text{Expected Risk}}.
\end{aligned}
\label{eq:soft_risk_clause}
\end{equation}

\subsubsection{Illustrative Example}
To concretize the proposed encoding, we consider the ADG substructure depicted in Fig. \ref{fig:ADG}.

\textbf{Hard Clauses:} Consider exploit node $e_1$, dependent on preconditions $c_1, c_2, c_4$. Eq. \eqref{eq:hard_exploit} encodes this AND-dependency as $C(e_1) = (\neg c_1 \vee \neg c_2 \vee \neg c_4 \vee e_1)$. This implies $e_1$ must be active if all three preconditions are present. For configuration $c_1$, mitigated by defense $d_1$, Eq. \eqref{eq:hard_defense} generates $C(c_1) = (\neg c_1 \vee \neg d_1) \wedge (c_1 \vee d_1)$.

\textbf{Soft Clauses:} For privilege node $p_1$, accessed via exploits $e_1$ or $e_2$, suppose the state vector is $\mathbf{v}=(1,0)$ (i.e., $e_1$ succeeds, $e_2$ fails). The blocking clause is $C_{\mathbf{v}} = (\neg e_1 \vee e_2)$ with weight $\rho_{e_1} \cdot I_{p_1}$. If the solver assigns Truth to $e_1$ and False to $e_2$, this clause evaluates to False, triggering the penalty.

The constructed WPM instance $\Phi$ encapsulates both the physical laws of the network (via $\Phi_{hard}$) and the strategic preferences of the administrator (via $\Phi_{soft}$). By solving $\Phi$, we implicitly perform a multi-objective search: the solver must strictly adhere to the attack graph topology while searching for a defense strategy $\mathbf{x}$ that  maximizes system utility.

\subsection{Theoretical Analysis}

% --- 过渡段：总结前文并引出证明 ---
The preceding sections have detailed the construction of the proposed framework, including the generation of hard clauses based on the causal dependencies in the Attack-Defense Graph and the formulation of soft clauses incorporating defense costs, potential security risks, configuration stability benefits, and attack success probabilities. 
To validate the mathematical rigor of our approach, we strictly prove the correctness of the proposed transformation. Specifically, the theoretical analysis focuses on three aspects. First, we prove that the WCNF encoding preserves the causal dependencies of the ADG. Second, we characterize how the encoded objective approximates the global utility $\mathcal{U}(\mathbf{x})$ through depth-wise probabilistic expansion. Third, we prove that the MaxSAT solver exactly optimizes the encoded objective.

\subsubsection{Preservation of Causal Dependencies}

To strictly demonstrate that the generated WCNF (Weighted Conjunctive Normal Form) instances accurately preserve the causal logic of the attack graph under specific defense strategies, we establish the following two lemmas regarding node state determination and then present the preservation theorem.

% --- 引理 1: 防御与配置的互斥性 ---
\begin{lemma}[Determinism of Configuration Nodes]
\label{lemma:config_determinism}
Given a defense strategy vector $\mathbf{x}$, the state of any configuration node $c \in N_p$ is uniquely determined and exhibits strict mutual exclusion with its corresponding defense action $d \in \mathbf{D}$.
\end{lemma}

\begin{proof}
Based on the hard clause construction in Eq. \eqref{eq:hard_defense}, a mutual exclusion constraint is defined as $d \leftrightarrow \neg c$. Consequently, determining $\mathbf{x}$ strictly fixes the truth values of all configuration nodes $N_p$, independent of other graph states.
\end{proof}

% --- 引理 2: 权限节点的完备性 ---
\begin{lemma}[Equivalence of Privilege Nodes]
\label{lemma:privilege_equivalence}
For any privilege node $p$, its truth value is strictly equivalent to the logical disjunction of its child exploit nodes. That is, $p \leftrightarrow \bigvee_{e \in child(p)} e$.
\end{lemma}

\begin{proof}
The causal logic for a privilege node implies that a privilege is gained if any contributing exploit is successful. The constructed constraints ensure that if any $e \in child(p)$ is True, $p$ is forced to True. Conversely, if all child exploits are False, no valid path exists to activate $p$. Thus, the state of $p$ satisfies the bidirectional equivalence $p \leftrightarrow \bigvee_{e \in child(p)} e$.
\end{proof}

\begin{theorem}[Preservation of Causal Logic]
\label{thm:preservation}
Given a fixed defense strategy $\mathbf{x}$, the WCNF optimum uniquely determines the system-configuration states, the privileges that the attacker may obtain, and their corresponding local compromise probabilities, even when the ADG contains directed cycles.
\end{theorem}

\noindent The proof is provided in Appendix~\ref{app:causal_proof}.

\subsubsection{Approximation to the Global Utility $\mathcal{U}(\mathbf{x})$}

We establish the approximation relationship between the encoded objective and the BP-based evaluation of the global utility $\mathcal{U}(\mathbf{x})$.

{\color{red}
\begin{theorem}[Depth-Wise Approximation of Global Utility]
\label{thm:depth_bp}
Equation~\eqref{eq:prob} provides a depth-one approximation of the compromise probabilities used in $\mathcal{U}(\mathbf{x})$. Each additional upstream exploit layer incorporates the corresponding BP messages into the encoded risk terms. When the expansion includes all messages used by the BP evaluation, the encoded compromise probabilities equal the BP beliefs, and the resulting objective coincides with the BP-based evaluation of $\mathcal{U}(\mathbf{x})$.
\end{theorem}

\noindent The proof is given in Appendix~\ref{app:bp_proof}. Deeper expansions reproduce more of the BP-based risk evaluation but increase the number of exploit-state combinations, and hence the number of risk clauses, exponentially. The current formulation therefore uses the depth-one approximation to balance global-utility fidelity and MaxSAT scalability.}

% --- 第二部分：证明软子句与目标函数的等价性 ---
\subsubsection{Equivalence of Objective Optimization}

Let $\mathcal{S}$ be the truth assignment space for all nodes in the ADG. The Weighted Partial MaxSAT solver implicitly optimizes a global objective function $\mathcal{F}(\mathbf{s})$, which serves as a computationally tractable approximation of the theoretical utility $\mathcal{U}(\mathbf{x})$ defined in Eq. \eqref{eq:global_function}. The solver searches for a state $\mathbf{s} \in \mathcal{S}$ that satisfies all topological constraints ($\Phi_{hard}$) while maximizing:

\begin{equation}
    \label{eq:solver_obj}
    \max_{\mathbf{s}}\, \mathcal{F}(\mathbf{s}) = \sum_{c =\text{True}} B_c - \sum_{d = \text{True}} C_d - \sum_{p = \text{True}} \Pr(p|\mathbf{s}) \cdot I_p,
\end{equation}
where the first two terms correspond to the aggregated utility of preserved configurations and the deployment cost of active defenses, respectively. The third term represents the cumulative penalty derived from risk mitigation failures.

It is important to clarify the notation used in Eq. \eqref{eq:solver_obj}:
\begin{itemize}
    \item \textbf{State Vector ($\mathbf{s}$ vs. $\mathbf{v}$):} Here, $\mathbf{s}$ denotes the global state vector encompassing the values of all node types ($c, p, e, d$) in the graph. This is distinct from the vector $\mathbf{v}$ introduced in the previous section, which represented the local subset of exploit states ($e \in child(p)$) specific to a single privilege node.
    
    \item \textbf{Local vs. Global Probability:} While the theoretical utility $\mathcal{U}(\mathbf{x})$ in Eq. \eqref{eq:global_function} depends on the global end-to-end attack probability $\Pr(p|\mathbf{x})$ to evaluate risk, computing this global reachability is computationally intractable. To bypass this, we utilize $\Pr(p|\mathbf{s})$, which captures the local compromise probability based strictly on the immediate state of $p$'s parents. By penalizing local attack combinations, $\mathcal{F}(\mathbf{s})$ provides a computationally tractable approximation for maximizing $\mathcal{U}(\mathbf{x})$, effectively suppressing overarching attack vectors through localized risk aggregation.
\end{itemize}

The Partial Weighted MAXSAT solver seeks an assignment that minimizes the sum of weights of \textit{falsified} soft clauses (often referred to as the cost of the solution).

\begin{theorem}[Optimization Equivalence]
\label{thm:optim}
Minimizing the total weight of falsified clauses in the MAXSAT formulation is mathematically equivalent to maximizing the global objective function $\mathcal{F}(\mathbf{s})$. Consequently, the assignment derived from the MAXSAT solver corresponds to the optimal assignment $\mathbf{s}^*$.
\end{theorem}

\begin{proof}
Based on the soft clause construction defined in Section \ref{subsec:transformation}, the MAXSAT solver minimizes the total weight of falsified clauses. This cost function, denoted as $\mathcal{K}_{sat}$, is explicitly expressed as:
\begin{equation}
    \min\, \mathcal{K}_{sat} = \sum_{d=\text{True}} C_d + \sum_{c=\text{False}} B_c + \sum_{p=\text{True}} \text{Pr}(p|s)I_p
\end{equation}
Substituting the complementary relationship for configuration weights, $\sum_{c=\text{False}} B_c = \sum_{c \in N_p} B_c - \sum_{c=\text{True}} B_c$, into the equation yields:
\begin{equation}
    \min\, \mathcal{K}_{sat} = W_{t} - \underbrace{\left( \sum_{c=\text{True}} B_c - \sum_{d=\text{True}} C_d - \sum_{p=\text{True}} \text{Pr}(p)I_p \right)}_{\mathcal{F}(\mathbf{s})},
\end{equation}
where $W_{t} = \sum_{c \in N_p}B_c$. 
Consequently, minimizing $\mathcal{K}_{sat}$ is strictly equivalent to maximizing the objective function $\mathcal{F}(\mathbf{s})$.
\end{proof}

\subsection{Resolution and Dynamic Adaptation}
\label{subsec:dynamic_adaptation}

\textcolor{red}{The WPM instance constructed in Section~\ref{subsec:transformation} represents the security posture at a static snapshot. However, network security involves evolving policies, observed adversarial actions, and newly hypothesized threats. Policy changes and observed events can be handled by updating constraints or objective terms, whereas newly hypothesized attack steps require extending the corresponding ADG scenario and regenerating the affected clauses. We therefore consider three forms of adaptation: policy-driven constraints, event-driven response, and uncertainty-driven scenario analysis.}

\textbf{Policy-Driven Constraints (Top-Down).} Administrators often possess domain knowledge or operational requirements that override automated optimization, typically manifesting as strict mandates on service availability or security zones. In our framework, these are handled by injecting \textit{Unit Hard Clauses} into $\Phi_{hard}$. Specifically, if a configuration node $c$ is deemed critical for business continuity, we append the unit clause $(c)$, forcing the solver to restrict the search space to solutions where $c=\text{True}$. Conversely, if a specific privilege $p$ is deemed unacceptable regardless of defense cost, we append the unit clause $(\neg p)$. 

For example, in the ADG shown in Fig. \ref{fig:ADG}, suppose the administrator determines that the network configuration represented by node $c_5$ is essential for business operations. We simply add the hard clause $(c_5)$ to the instance. Consequently, the solver is mathematically barred from selecting any defense strategy that would negate $c_5$. This constraint forces the optimization engine to seek alternative hardening paths, such as deploying other defenses or blocking upstream exploits, to break the attack chain leading to the target privilege $p_1$, ensuring security without compromising the mandatory business utility.

\textbf{Event-Driven Response (Bottom-Up).} When the network operates in real-time, an Intrusion Detection System (IDS) may alert that a specific exploit $e_{obs}$ has been successfully executed. According to the topological logic defined in $\Phi_{hard}$ (specifically Eq. \eqref{eq:hard_privilege}), the occurrence of $e_{obs}$ deterministically implies the compromise of the consequent privilege $p$ (i.e., $p=\text{True}$). This transition marks the shift from a ``probabilistic risk" to a ``realized loss." 

In this scenario, the dynamic adaptation requires no modification to the structural constraints ($\Phi_{hard}$ remains invariant). Instead, we perform an \textit{Objective Adjustment}. The original soft clauses penalizing the occurrence of $e_{obs}$ (or the specific state vector $\mathbf{v}$ leading to $p$) now represent a sunk cost. To reflect this, we explicitly remove these corresponding clauses from $\Phi_{soft}$. By eliminating the penalty associated with the already-compromised node, the solver automatically shifts its optimization focus. It ceases allocating resources to prevent the irreversible event $e_{obs}$ and instead searches for the optimal strategy to minimize the residual risk, prioritizing defenses against subsequent attack steps that depend on the compromised privilege $p$.

\textcolor{red}{\textbf{Uncertainty-Driven Scenario Analysis.} Potential zero-day attacks are treated as hypothetical scenarios rather than known CVE-based vulnerabilities. Using the existing exploit-node semantics, a scenario can include an additional exploit node connected to its assumed prerequisites and resulting privilege. Because no CVSS score is available, its success probability is specified as a scenario prior informed by historical attack observations, threat intelligence, or expert knowledge. Alternative assumptions can be examined by varying this prior. When threat intelligence indicates a previously unmodeled attack path, the scenario-specific ADG is extended with the corresponding node and dependencies, and the affected causal and risk clauses are regenerated.}

In summary, the proposed framework establishes a rigorous computational paradigm that transforms network hardening from heuristic graph analysis into a Weighted Partial MaxSAT optimization task. By constructing a strict isomorphism between graph topology and Boolean constraints, this approach preserves the semantic fidelity of attack-defense dependencies while unlocking the computational efficiency of modern solvers. Crucially, our formulation bridges the gap between probabilistic risk assessment and discrete logical reasoning by encoding expected attack impacts into weighted soft clauses, enabling the optimizer to account for multi-step attack likelihoods within the trade-off space. Coupled with the mechanism for incremental constraint adaptation, this methodology effectively converts static security planning into a dynamic resilience engine, capable of delivering near-optimal, operationally viable strategies in response to real-time adversarial evolution.

\section{Evaluation}
\label{sec:experiments}

In this section, we present a comprehensive empirical evaluation of the proposed logic-based optimization framework. \textcolor{red}{Our experiments are designed to rigorously assess four critical dimensions of the system's performance: (1) the efficacy of the high-quality strategies derived under complex security-economic trade-offs; (2) the adaptability of the model to evolving administrative policies and real-time intrusion events; (3) the comparative performance of the framework against representative optimization baselines; and (4) the computational feasibility of the approach when deployed in large-scale network environments.}

\subsection{Experimental Setup}
\label{subsec:setup}

In this section, we present the experimental evaluation of the proposed framework. We detail the construction of a representative enterprise network testbed \textcolor{red}{implemented using virtual machines}, the instantiation of the Attack-Defense Graph based on real-world vulnerability data, and the configuration of the optimization environment. The experimental design aims to evaluate the framework's capability in generating a high-quality defense strategy under stochastic attack conditions.

\begin{figure}[htbp]
\centering
\includegraphics[width=\linewidth]{testnetwork.pdf}
\caption{Topology of the experimental enterprise testbed. The network is segmented into three logical zones (DMZ, Trusted, Inside) to simulate a defense-in-depth architecture.}
\label{fig:testnetwork}
\end{figure}

\textbf{Testbed Topology and Zoning:} 
To rigorously evaluate the proposed framework, we deployed a segmented network testbed modeled after a hierarchical enterprise environment. As illustrated in Fig. \ref{fig:testnetwork}, the network is strictly partitioned into three logical security zones to enforce defense-in-depth principles. 
The Demilitarized Zone (DMZ) serves as the presentation layer, hosting public-facing services including a Web Server, an API Gateway, a Mail Server, and a DNS Server. 
The Trusted Zone functions as the application and data layer, hosting critical infrastructure components such as the Core Database and a File Server. Notably, an Edge Gateway is also deployed in this zone to facilitate remote SSH administration. 
The innermost layer is the Inside Zone (also referred to as the Restricted Management Zone), which hosts the high-value Admin Workstation. This workstation represents the target asset, as it holds privileged credentials for managing the entire infrastructure.

\textbf{Network Accessibility:} 
Both the DMZ services (Web, Mail, API, DNS) and the Edge Gateway are directly accessible from the Internet. 
Internal traffic flows are restricted by functional requirements: compromised DMZ nodes can access specific resources in the Trusted Zone to support backend operations; for instance, the Web Server interacts with the Core Database, while the Mail Server utilizes the File Server for storage. 
Crucially, the Trusted Zone retains administrative connectivity to the Admin Workstation in the Inside Zone to facilitate remote management. 

\textbf{Vulnerability Profiling:} 
We utilized Tenable Nessus Essentials~\cite{yoran2002nessus} to perform a comprehensive vulnerability scan of the deployed services. The identified vulnerabilities were mapped to the National Vulnerability Database to extract Common Vulnerabilities and Exposures (CVE) entries. 
In the generated Attack-Defense Graph, these CVE entries are represented as configuration nodes (preconditions). 
To quantify the stochastic nature of vulnerability exploitation, the CVSS v3 Exploitability Score of each vulnerability is scaled down by a factor of 10 to derive the probability parameter $\rho_{e_i} \in (0,1]$~\cite{FIRST2015-CVSSv30, munoz2017exact}. This parameter defines the intrinsic success probability of the exploit node $e_i$, directly reflecting the ease and technical means required to compromise the underlying vulnerable component. 
For exploit nodes that do not rely on specific software vulnerabilities (e.g., logical lateral movement), we assume deterministic execution given the preconditions are met, setting $\rho_{e_i}=1$.
Table \ref{tab:vuln_profile} summarizes the critical CVEs acting as preconditions and the derived success probabilities for their associated exploits.

\begin{table}[htbp]
\caption{Vulnerability Configurations and Associated Exploit Probabilities}
\label{tab:vuln_profile}
\centering
\renewcommand{\arraystretch}{1.3} % 保持舒适的行高
% 使用 resizebox 强制表格适应单栏宽度
\resizebox{\columnwidth}{!}{%
\begin{tabular}{|l|l|c|}
\hline
\textbf{Host Node} & \textbf{CVE Configuration} & \textbf{Exploit Prob. ($\rho_{e_i}$)} \\
\hline \hline
Web Server (DMZ) & CVE-2021-44228 & 0.39 \\ \hline
API Server (DMZ) & CVE-2021-41773 & 0.39 \\ \hline
Mail Server (DMZ) &\color{red}{ CVE-2019-15846} &\color{red}{ 0.39 }\\ \hline
Edge Gateway (Trusted) & CVE-2018-15473 & 0.39 \\ \hline
Core Database (Trusted) & CVE-2019-9193 & 0.12 \\ \hline
File Server (Trusted) & CVE-2017-7494 & 0.39 \\ \hline
\multirow{2}{*}{Admin Workstation (Inside)} & \color{red}{CVE-2019-1181} & 0.39 \\ \cline{2-3}
                                            & CVE-2020-0796 & 0.39 \\ \hline
\end{tabular}%
}
\end{table}

\textbf{Candidate Defense Strategies:} 
We defined a repository of defense countermeasures with varying granularities to challenge the solver's decision-making. As shown in Table \ref{tab:defense_plan}, these strategies range from surgical remediation (e.g., $D_{1A}, D_7$), which strictly removes the vulnerability node while preserving service availability, to drastic isolation (e.g., $D_{10}, D_{11}$), which deactivates the service entirely by removing associated connectivity rules (\texttt{hacl}) and service definitions (\texttt{networkServiceInfo}) alongside the vulnerability.

\begin{table}[htbp]
\caption{Defense Countermeasures and Operational Scope}
\label{tab:defense_plan}
\centering
\renewcommand{\arraystretch}{1.2}
\resizebox{\columnwidth}{!}{%
\begin{tabular}{|l|l|l|}
\hline
\textbf{Target Host} & \textbf{ID} & \textbf{Defense Mechanism} \\ \hline \hline

% --- Admin Workstation ---
\multirow{4}{*}{\textbf{Admin Workstation}} 
 & $D_{1A/B}$ & Targeted Patching (RDP/SMB Service) \\ \cline{2-3}
 & $D_{2}$ & Service Hardening (Config Tuning) \\ \cline{2-3}
 & $D_{3E/I}$ & Inbound Traffic Filtering (ACL) \\ \cline{2-3}
 & $D_{4}$ & Internal Segmentation (Block Lateral) \\ \hline

% --- Edge Gateway ---
\multirow{2}{*}{\textbf{Edge Gateway}} 
 & $D_{5V}$ & Targeted Patching (SSH Service) \\ \cline{2-3}
 & $D_{5C/A}$ & Auth Enforcement / Source Restriction \\ \hline

% --- Core Database ---
\multirow{2}{*}{\textbf{Core Database}} 
 & $D_{8V/C}$ & Patching / Service Hardening \\ \cline{2-3}
 & $D_{9}$ & Application-Layer Access Control \\ \hline

% --- DMZ Services (Key Logic Change) ---
\multirow{4}{*}{\textbf{DMZ Services}} 
 & $D_{6}$ & Targeted Patching (API Gateway) \\ \cline{2-3}
 & $D_{7}$ & Targeted Patching (Web Server) \\ \cline{2-3}
 & $D_{10}$ & \textbf{Service Deactivation} (Disable File Srv \& SMB) \\ \cline{2-3}
 & $D_{11}$ & \textbf{Service Deactivation} (Disable Mail Srv \& SMTP) \\ \hline

\end{tabular}%
}
\end{table}

\textbf{Implementation and Computational Environment:} 
The proposed optimization framework is implemented in Python, integrating the MaxHS~\cite{davies2011solving} solver to handle the weighted Partial MaxSAT formulation. MaxHS was selected for its efficiency in handling structured dependencies via the Implicit Hitting Set (IHS) approach. To ensure rigorous evaluation, all experiments were conducted on a Linux workstation equipped with an AMD Ryzen Threadripper PRO 9955WX processor (running at 4.5~GHz) and 188~GB of RAM. The solver execution was restricted to a single thread to ensure consistent and precise CPU timing measurements.

\textbf{Baselines and Evaluation Metrics:}
\textcolor{red}{We compare MaxSAT with the Bayesian attack-graph GA of Poolsappasit et al.~\cite{poolsappasit2011dynamic}, the probabilistic min-max MILP of Khouzani et al.~\cite{KHOUZANI2019894}, and the monotonic-gradient method of Zenitani~\cite{zenitani2022multi}. All methods use the same attack graphs, parameters, and candidate defenses, and their returned portfolios are evaluated using the common utility $\mathcal{U}$. We use exact variable elimination on small graphs and loopy belief propagation (100 iterations, damping 0.2, and tolerance $10^{-6}$) on medium graphs. The GA follows its original configuration, and we report the median utility and runtime over five fixed-seed runs per instance. Khouzani-MILP minimizes worst-path risk across ten uniformly spaced defense budgets, from which we select the portfolio with the highest $\mathcal{U}$. Zenitani follows its original search procedure, using a resource-bounded medium-scale configuration with five refinement rounds, five neighbor samples, and descent caps of 10 and 40 for structured and random graphs, respectively. Runtime is measured from the in-memory attack graph and parameter table to the final defense portfolio. MaxSAT exactly optimizes the encoded surrogate $\mathcal{F}$, whereas all comparisons use the common ex-post utility $\mathcal{U}$. Throughout all experiments, two methods are considered tied only when their final evaluated utility values are equal within floating-point precision.}

\subsection{Effectiveness Analysis: A Case Study}
\label{subsec:case_study}

To validate the decision-theoretic capabilities of the proposed framework, we analyze a critical multi-stage attack scenario extracted from the experimental environment. The ultimate goal of our modeling framework is to determine the optimal defense strategy that maximizes system utility. This objective requires balancing three competing dimensions: preserving System Operational Benefits, minimizing Defense Implementation Costs, and mitigating Security Risks.

\subsubsection{Attack-Defense Graph Construction}

\begin{figure}[htbp]
\centering
% Placeholder for the ADG Fig
\includegraphics[width=\linewidth]{ADG_1.pdf}
\caption{Visualization of the Attack-Defense Graph focusing on the critical path to the Admin Workstation. Defense nodes (green) act as intervention points to break the attack chain (red).}
\label{fig:adg_case_study}
\end{figure}

The experimental environment is modeled using a high-fidelity Attack-Defense Graph, generated through a three-stage automated pipeline. 
First, Tenable Nessus is deployed to perform comprehensive vulnerability profiling, extracting CVE entries across the segmented network. 
Second, these vulnerability signatures, combined with network topology and host configuration rules, are ingested into the MulVAL reasoning engine. This step constructs the logical attack graph, mapping causal dependencies where atomic exploits are chained to reach deep network assets. 
Finally, the framework augments this baseline graph by injecting a layer of candidate defense nodes. These nodes represent actionable interventions, such as targeted patching and network segmentation, that can logically negate specific preconditions in the graph. \textcolor{red}{The resulting ADG, visualized in Fig.~\ref{fig:adg_case_study}, provides a causal representation of the attack surface captured by the supplied network information and attack-rule base.}

\subsubsection{Optimal Defense Selection and Strategic Rationale}
\label{subsec:optimization_analysis}

To evaluate the framework's decision-theoretic efficacy, we formulated the defense selection as a utility maximization problem. This involves a rigorous trade-off among three dimensions: preserving Operational Benefits, minimizing Defense Implementation Costs, and mitigating Security Risks. In our model, we differentiated between ``Service Utility" (high value, derived from core services like SQL or RDP) and ``Connectivity Utility" (low value, derived from routing paths). This stratification ensures that the optimization model inherently prioritizes the continuity of critical applications, treating network routes as sacrificial edges only when necessary to prevent catastrophic loss.

Leveraging the MaxHS solver, the framework computed the globally optimal defense strategy within 0.005 seconds. The resulting active defense set is:
\begin{equation}
    D_{active} = \{ D_{3E},D_4, D_{5A}, D_9 \}.
\end{equation}
It is worth noting that this result, derived through our logic-based discrete optimization, aligns perfectly with solutions obtained using precise probabilistic risk assessment models. This convergence validates that our discretized cost-utility parameters accurately capture the stochastic nature of the threat landscape without the computational overhead of continuous probability calculations.

The rationale behind this outcome reveals a sophisticated risk-reward calculus, particularly when contrasted with a deterministic baseline scenario where vulnerability exploitability rates are disregarded (i.e., assuming $\rho = 1$). In such a deterministic baseline, the optimization yields a markedly different set: $D_{baseline} = \{  D_{3E}, D_{5A}, D_9, D_6, D_{11} \}$.
\begin{itemize}
    \item The deterministic model, operating under the assumption of guaranteed compromise, enforces drastic measures on the perimeter. It deploys the high-cost patch for the API Gateway ($D_6$) and triggers Service Deactivation for the Mail Server. While this strategy strictly secures the perimeter, it is suboptimal for two reasons: it incurs prohibitive operational friction (Defense Cost) and, more critically, it significantly degrades network configuration benefits by disabling the valid Mail Service configuration, resulting in a substantial loss of Service Utility.
    \item In contrast, our proposed model identifies that the expected loss from these peripheral nodes is low (due to $\rho < 1$). Consequently, it preserves the high-value configurations of the Mail and API services (avoiding the utility loss of $D_{11}$ and cost of $D_6$). Instead, it relies on the low-cost segmentation rule $D_4$ (Block Lateral to Admin). This approach secures high-value targets while severing only non-essential communication links, thus maximizing net utility without significantly degrading routine system functions.
\end{itemize}

This comparison highlights the framework's ability to implement ``Smart Segmentation", replacing brute-force service disruption with cost-effective internal containment. Thus, the system correctly distinguishes between existential threats that must be blocked and statistical anomalies that can be tolerated for the sake of operational continuity.

\subsection{Dynamic Adaptability Evaluation}
\label{subsec:dynamic_flexibility}

While static optimization establishes a foundational baseline, a robust defense framework must demonstrate resilience against evolving administrative mandates and real-time intrusion events. A critical advantage of our logic-based formalism is its inherent plasticity, which allows the system to dynamically recalibrate the defense posture by injecting unit clauses or modifying constraints without the computational overhead of restructuring the underlying graph topology. We evaluate this capability through two distinct operational scenarios using the proposed framework.

The first scenario, Policy-Driven Adaptation, examines the system's response to rigid administrative constraints. In the unconstrained optimization (Section \ref{subsec:optimization_analysis}), the solver prioritized the segmentation strategy $D_4$ (Block Lateral Movement) to neutralize the risk to the Admin Workstation. However, operational requirements often dictate that specific data flows must remain open. To simulate this, we injected a hard constraint enforcing the persistence of configuration node $C_{20}$ (\texttt{hacl(db\_core, admin\_ws, tcp, 445)}), which represents a critical SMB channel for logs. Upon re-optimization with this constraint, the solver automatically pivoted its strategy. Since the low-cost segmentation option $D_4$ was rendered infeasible by the mandatory retention of $C_{20}$, the framework converged on a new optimal set: $\{ D_{1B}, D_{3E}, D_{5A}, D_9 \}$. Crucially, the system recommended the high-friction remediation strategy $D_{1B}$ (Patching the Admin Workstation's SMB service). This shift confirms that the system can strictly adhere to hierarchical business continuity policies, correctly identifying that minimizing net system loss now requires absorbing the high cost of patching to neutralize the existential risk, as the cheaper network isolation option is policy-blocked.

The second scenario, Event-Driven Response, evaluates the framework's agility in critical post-breach situations. We simulated a deep intrusion event where the IDS confirms that the attacker has bypassed the perimeter and acquired Network Access to the Admin Workstation (represented by asserting $P_2$: \texttt{netAccess(admin\_ws, tcp, 3389)} as True), a state representing an immediate existential threat as the adversary is merely one exploit step away from root compromise ($P_1$). In response, the solver's strategy underwent a fundamental shift, generating the active defense set $\{ D_{1A},  D_{3E}, D_4, D_{5A}, D_9 \}$. Most notably, the solver activated $D_{1A}$ (Patching the Admin Workstation's RDP service), a strategy it had previously rejected due to its prohibitive operational friction. This decision illustrates the concept of ``Emergency Hardening''. In the pre-breach state, the probabilistic reliance on segmentation ($D_4$) was sufficient and cost-effective; however, once the attacker successfully acquired $P_2$, the utility of blocking the inbound path diminished as the perimeter was already breached. Consequently, the probability of $P_1$ compromise surged drastically, now bounded only by the intrinsic exploit success rate. Faced with this elevated risk, the solver prioritized the high-cost patch ($D_{1A}$) as an essential defense to prevent full compromise. This behavior underscores the framework's capacity for dynamic adaptation by explicitly trading operational efficiency for system resilience when critical security thresholds are breached.

\subsection{Comparative and Scalability Evaluation}

{\color{red}
\subsubsection{Synthetic Datasets and Topology Specifications}

To support comparative evaluation and scalability testing, we generate two complementary synthetic graph families. To maintain structural plausibility, both families preserve the typed semantics of privilege ($P$), exploit ($E$), configuration ($C$), and defense ($D$) nodes. Under the dependency-edge convention of the ADG, the admissible relations are $P\rightarrow E$, $E\rightarrow(C\cup P)$, and $C\rightarrow D$. Both families also enforce acyclicity and bounded local connectivity, and edges are sampled only between semantically admissible node types. In addition, at most five immediate exploit alternatives may lead to a privilege~\cite{munoz2017exact}, reflecting the finite local dependencies considered in prior Bayesian attack graph evaluations.

\begin{itemize}
    \item \textbf{Structured Hierarchical Graphs ($G_s$):}
    The generator constructs a rooted privilege hierarchy in which every privilege is reachable from the entry privilege and privilege transitions follow a topological order. Exploit, configuration, and defense nodes are then added according to the prescribed dependency pattern. Each exploit is associated with one to three configurations, each configuration is shared by at most five exploits, and each defense covers at most three configurations. This construction captures multi-stage attack progression, alternative attack paths, and the reuse of configurations and defenses in hierarchically organized networks.

    \item \textbf{Unstructured Random Graphs ($G_r$):}
    Serving as a theoretical stress-test baseline, $G_r$ abandons the fixed hierarchical backbone of $G_s$ in favor of stochastic cross-layer mappings. The generator samples semantically admissible endpoints while retaining the degree and acyclicity constraints. Every privilege participates in at least one exploit dependency, and privilege transitions follow a topological order. This construction creates irregular branching, overlapping attack paths, and shared prerequisites, providing a structurally valid but computationally demanding workload for scalability evaluation.
\end{itemize}

Together, $G_s$ and $G_r$ cover hierarchical and irregular dependency patterns while preserving typed causal relations, sparse local dependencies, and multi-stage attack progression. These constraints allow controlled variation in graph size and topology without discarding the principal structural characteristics of real attack-defense graphs.



\subsubsection{Comparison with Optimization Baselines}
We benchmark the four methods on small- and medium-scale instances. The ADG size ranges from 50 to 100 nodes for small instances and from 700 to 1,200 nodes for medium instances.

\begin{figure}[htbp]
\centering
\includegraphics[width=\linewidth]{baselines_comparison.pdf}
\caption{
Pairwise comparison of defense quality and runtime.
The left panel reports the fraction of instances in which the baseline attains higher utility, the two methods are tied, or MaxSAT attains higher utility.
The right panel reports the mean end-to-end runtime on a logarithmic scale.
For GA, the per-instance quality and runtime are represented by the median over five runs.
Each small dataset contains 100 instances, whereas each medium dataset contains 10 instances.
}
\label{fig:baselines_comparison}
\end{figure}

Figure~\ref{fig:baselines_comparison} shows that MaxSAT achieves competitive solution quality on small graphs. On structured instances, MaxSAT matches or exceeds GA, Khouzani-MILP, and Zenitani in 69\%, 74\%, and 69\% of cases, respectively; the corresponding proportions on random instances are 92\%, 96\%, and 92\%. The advantage becomes more pronounced on medium graphs. MaxSAT achieves higher evaluated utility than GA on six of the ten structured instances and than both Khouzani-MILP and Zenitani on all ten. On random medium graphs, MaxSAT achieves higher utility than every baseline on all ten instances.

The runtime results further demonstrate the scalability advantage of the proposed formulation. MaxSAT has the lowest mean end-to-end runtime in all four datasets. On medium structured graphs, MaxSAT requires 0.038 seconds on average, compared with 3,443.1 seconds for GA, 0.334 seconds for Khouzani-MILP, and 873.1 seconds for Zenitani. On medium random graphs, the corresponding runtimes are 10.1, 545.1, 136.0, and 1,098.8 seconds.

The results reflect differences in both computational workflow and optimization focus. GA and Zenitani require repeated portfolio evaluations during evolutionary search and neighborhood refinement, which increasingly limits their scalability as the graph and defense spaces grow. While Khouzani scales well, its objective emphasizes the worst attack path, providing a less comprehensive assessment of security benefits, expected losses, and defense costs across the network. By jointly encoding feasibility and the local risk surrogate in a single WCNF instance, MaxSAT exploits the underlying logical structure to obtain defense portfolios with higher evaluated global utility at low computational cost.

\subsubsection{Local-Degree Stress Evaluation}

According to Eq.~\eqref{eq:risk_clause}, a privilege node with \(k\) direct exploit children generates \(2^k-1\) risk clauses. We evaluate \(k_{\max}\in\{5,7,10\}\) on separate medium-scale datasets, each containing ten structured and ten random graphs. All three settings use newly generated datasets under the same protocol. Degrees follow truncated discrete-normal distributions over \([1,k_{\max}]\), with \((\mu,\sigma)=(3.0,0.8)\), \((4.0,1.2)\), and \((5.5,1.8)\), respectively. W/T/L summarizes MaxSAT's 30 pairwise comparisons against the three baselines under \(\mathcal{U}\) for each setting.

\begin{table}[t]
\centering
\caption{Local-degree stress results.}
\label{tab:k_stress}
\footnotesize
\setlength{\tabcolsep}{3pt}
\renewcommand{\arraystretch}{1.1}
\begin{tabular}{llrrr}
\toprule
\textbf{Topology} &
\(\boldsymbol{k_{\max}}\) &
\textbf{Risk Clauses} &
\textbf{MaxSAT (s)} &
\(\boldsymbol{\mathcal{U}}\) \textbf{W/T/L} \\
\midrule
Structured & 5  & 821   & 0.036   & 29/1/0 \\
Structured & 7  & 2,153 & 0.079   & 30/0/0 \\
Structured & 10 & 9,299 & 0.781   & 30/0/0 \\
Random     & 5  & 828   & 1.852   & 30/0/0 \\
Random     & 7  & 2,078 & 9.762   & 30/0/0 \\
Random     & 10 & 9,170 & 203.904 & 30/0/0 \\
\bottomrule
\end{tabular}
\end{table}

As shown in Table~\ref{tab:k_stress}, clause count and runtime increase with \(k_{\max}\), while MaxSAT retains the lowest mean runtime among the four methods in every setting. Using full-graph \(\mathcal{U}\) evaluated by loopy belief propagation, MaxSAT records 179 wins, one tie, and no losses, showing that the larger tested degrees increase computational cost without reducing defense quality.

\subsubsection{Comparative Parameter Sensitivity}
\label{subsec:parameter_sensitivity}

We evaluate whether conclusions about comparative solution quality are sensitive to uncertainty in the externally specified economic parameters. The experiment uses 20 medium graphs, comprising ten structured and ten random instances. We perturb one parameter class at a time: configuration benefit \(B_c\), potential loss from privilege compromise \(I_p\), or defense implementation cost \(C_d\). For each applicable node, the relative perturbation magnitude is sampled independently and uniformly between zero and the specified upper bound, and the parameter is randomly increased or decreased. For the same node, the direction and relative perturbation pattern are retained across the \(10\%\), \(20\%\), and \(30\%\) settings, with the perturbation magnitude scaled proportionally to the corresponding upper bound. Graph topology, candidate defenses, exploit probabilities, and the other parameter classes remain unchanged. This procedure produces 180 perturbed instances. MaxSAT, GA, Khouzani, and Zenitani are applied to the same instances and evaluated under the perturbed parameters, following the common evaluation protocol described in Section~\ref{subsec:setup}.

From the perspective of MaxSAT, the W/T/L counts against GA are 51/3/6, 50/4/6, and 49/6/5 for the \(10\%\), \(20\%\), and \(30\%\) perturbation ranges, respectively, yielding 150 wins, 13 ties, and 17 losses overall. MaxSAT wins all 180 comparisons against both Khouzani and Zenitani. Across topology classes, MaxSAT records 90/0/0 against GA on random graphs and 60/13/17 on structured graphs. The corresponding numbers of wins or ties against GA are 54, 54, and 55 as the perturbation range increases. Overall, MaxSAT remains predominantly superior to GA and uniformly superior to Khouzani and Zenitani across all perturbation levels, supporting its robust comparative advantage within the tested ranges.
}

\subsubsection{Large-Scale Scalability Stress Test}

To further evaluate the scalability of the MaxSAT framework, we perform stress tests using large-scale synthetic datasets. For structured topologies, we deployed the solver on 100 dense instances, each exceeding 15,000 variables and 40,000 clauses. Remarkably, MaxHS consistently converged to the exact optimal solution of the approximated objective $\mathcal{F}(\mathbf{s})$ in less than 1 second per instance. This performance validates that the solver effectively exploits the hierarchical dependencies within structured networks to natively decompose the search space.

We next consider random topologies as a stress-test setting.
Compared with structured graphs, random instances contain more entangled dependencies and weaker modularity, making them harder for the solver to decompose.
Fig.~\ref{fig:random_convergence} reports the convergence of the MaxSAT objective $\mathcal{F}(\mathbf{s})$ over wall-clock time on random instances ranging from $7{,}900$ to $15{,}800$ variables and from $23{,}712$ to $48{,}478$ clauses. Across all tested scales, MaxHS improves the incumbent solution rapidly in the early solving stage.
For the largest instance, with $15{,}800$ variables and $48{,}478$ clauses, the objective value reaches a stable plateau within approximately two minutes.
The subsequent marginal improvement indicates that the solver has already obtained a near-final solution despite the lack of exploitable structure.
These results show that the proposed WCNF encoding remains tractable under substantially enlarged and structurally challenging instances.

Overall, the framework scales well on both topology classes: structured enterprise-like graphs are solved within seconds, while random graphs still exhibit fast objective convergence at larger scales.
This confirms the practical scalability of the proposed MaxSAT-based hardening formulation.

\begin{figure}[t]
    \centering
    \includegraphics[width=\linewidth]{convergence_objective_v2.pdf}
    \caption{
    Convergence of the MaxSAT objective $\mathcal{F}(\mathbf{s})$ on large-scale random instances.
    The instances range from $7{,}900$ to $15{,}800$ variables and from $23{,}712$ to $48{,}478$ clauses.
    MaxHS rapidly improves the incumbent objective value and reaches a stable plateau across all tested scales, indicating strong scalability under random dependency structures.
    }
    \label{fig:random_convergence}
\end{figure}


\section{Conclusion}
\label{sec:conclusion}

{\color{red}
In this paper, we presented a framework based on Weighted Partial MaxSAT for network defense hardening that integrates attack causality, exploitation likelihoods, and defense countermeasure costs into a fine-grained analysis. By mapping defense nodes to network configurations, we constructed an Attack-Defense Graph and transformed the hardening problem into WCNF. Hard clauses encode causal dependencies and defense mappings, while weighted soft clauses capture preserved operational benefit, aggregate deployment cost, and risk penalties derived from local conditional probabilities. Theoretical analysis establishes that the encoding preserves causal dependencies and that MaxSAT exactly optimizes the encoded surrogate objective $\mathcal{F}(\mathbf{s})$. It further characterizes how the surrogate relates to the BP evaluation of global utility $\mathcal{U}(\mathbf{x})$ as the expansion depth increases.

Under a common protocol for evaluating global utility, MaxSAT achieves competitive solution quality on small graphs and increasingly clear quality and runtime advantages on medium graphs. It obtains the lowest mean runtime on all four comparative datasets, while the experiments on local degree and parameter sensitivity demonstrate robustness over the tested ranges. The case studies further illustrate adaptation to policy constraints and observed attack events.

These results remain conditional on the completeness of the supplied ADG and the accuracy of its parameters. The current formulation uses externally specified benefits, losses, and aggregate defense costs, assumes binary defense effects, and adopts a risk approximation at depth one that generates $2^k-1$ clauses per privilege node. Future work will improve graph updating and parameter calibration, distinguish cost components, support partial and combined defense effects, and investigate adaptive probability expansion.
}



% conference papers do not normally have an appendix


% use section* for acknowledgment
%\section*{Acknowledgment}


%The authors would like to thank...





% trigger a \newpage just before the given reference
% number - used to balance the columns on the last page
% adjust value as needed - may need to be readjusted if
% the document is modified later
%\IEEEtriggeratref{8}
% The "triggered" command can be changed if desired:
%\IEEEtriggercmd{\enlargethispage{-5in}}

% references section

% can use a bibliography generated by BibTeX as a .bbl file
% BibTeX documentation can be easily obtained at:
% http://mirror.ctan.org/biblio/bibtex/contrib/doc/
% The IEEEtran BibTeX style support page is at:
% http://www.michaelshell.org/tex/ieeetran/bibtex/
%\bibliographystyle{IEEEtran}
% argument is your BibTeX string definitions and bibliography database(s)
%\bibliography{IEEEabrv,../bib/paper}
%
% <OR> manually copy in the resultant .bbl file
% set second argument of \begin to the number of references
% (used to reserve space for the reference number labels box)
\bibliographystyle{IEEEtran}
\bibliography{refs}

\appendix

\subsection{Proof of Theorem~\ref{thm:preservation}}
\label{app:causal_proof}

Theorem~\ref{thm:preservation} establishes that, given a fixed defense strategy, the WCNF optimum uniquely determines the system-configuration states, the privileges that the attacker may obtain, and their corresponding local compromise probabilities, even when the ADG contains directed cycles.

\begin{proof}
By Lemma~\ref{lemma:config_determinism}, the fixed defense strategy $\mathbf{x}$ uniquely determines all configuration-node states. By Lemma~\ref{lemma:privilege_equivalence}, a privilege node is True if and only if at least one of its child exploit nodes is True. When the dependencies are acyclic, the node states follow directly from these constraints. When they contain a directed cycle, the hard clauses may admit more than one assignment, so the effect of the soft clauses must also be considered.

Consider a directed dependency cycle

$$
p_1\rightarrow e_1\rightarrow p_2\rightarrow e_2\rightarrow\cdots\rightarrow p_m\rightarrow e_m\rightarrow p_1,
$$
where each arrow points from a node to one of its dependencies and all configuration-node preconditions of the exploit nodes are True. Without external privilege support, the hard clauses admit both the all-False and all-True assignments: the former satisfies each exploit clause through a False cyclic privilege precondition, whereas the latter satisfies each exploit clause through its True exploit literal and each privilege clause through a True child exploit.

The soft clauses exclude the unsupported all-True assignment. It incurs a positive risk penalty for every privilege in the cycle, whereas the all-False assignment incurs no risk penalty and leaves the benefit and defense-cost terms unchanged. Therefore, only the all-False assignment can be optimal. More generally, an unsupported exploit that is not forced by its preconditions is assigned False whenever setting it to True increases the local risk penalty.

If a privilege in the cycle is supported by an attack path outside the cycle, the corresponding exploit and privilege states are forced through Eqs.~\eqref{eq:hard_exploit} and~\eqref{eq:hard_privilege} whenever their remaining preconditions are satisfied. Thus, the optimal assignment retains externally supported privileges and excludes privileges supported only by the cycle.

Equation~\eqref{eq:prob} is monotonically nondecreasing in every exploit state, so changing $v_i$ from 0 to 1 can only increase or preserve the local compromise probability. Equation~(14) constructs the corresponding blocking clause whose violated weight, $\Pr(p\mid\mathbf{v})I_p$, represents the resulting risk loss. Therefore, an unforced exploit is assigned False whenever its activation increases the risk loss. If the risk loss remains unchanged, the exploit assignment may differ, but the local probability is unchanged. Consequently, all optimal assignments determine the same system configurations, obtainable privileges, and corresponding local probabilities.
\end{proof}

\subsection{Depth-Wise Probability Expansion and BP Correspondence}
\label{app:bp_proof}

Let $q_P^{(d)}$ denote the compromise probability of a target privilege node $P$ when its upstream exploit dependencies are expanded to a maximum depth $d$. Different branches may terminate at different depths. We therefore use $m_z$ to denote the probability message propagated from an internal or boundary node $z$ toward the target privilege.

Consider a privilege node $P_1$ directly reachable through exploits $E_1$ and $E_2$. At depth one, both exploits are boundary nodes with messages
\begin{equation}
m_{E_1}=\rho_1v_1,
\qquad
m_{E_2}=\rho_2v_2.
\end{equation}
The compromise probability of $P_1$ is
\begin{equation}
\begin{aligned}
q_{P_1}^{(1)}
&=1-(1-m_{E_1})(1-m_{E_2})\\
&=1-(1-\rho_1)^{v_1}(1-\rho_2)^{v_2},
\end{aligned}
\label{eq:depth_one_example}
\end{equation}
which is Eq.~\eqref{eq:prob}.

Suppose that $E_1$ requires privilege $P_2$, which can be independently obtained through exploits $E_3$, $E_4$, and $E_5$. Expanding the $E_1$ branch introduces the boundary messages
\begin{equation}
m_{E_i}=\rho_iv_i,
\qquad i\in\{3,4,5\}.
\end{equation}
These messages are aggregated at $P_2$ as
\begin{equation}
m_{P_2}
=
1-\prod_{i=3}^{5}(1-m_{E_i}).
\label{eq:p2_message}
\end{equation}
Assuming that the non-privilege prerequisites of $E_1$ are satisfied, the message propagated through $E_1$ becomes
\begin{equation}
m_{E_1}
=
\rho_1m_{P_2}.
\label{eq:e1_message}
\end{equation}
Since the $E_2$ branch remains unexpanded, its message remains $m_{E_2}=\rho_2v_2$. The resulting depth-two probability is
\begin{equation}
q_{P_1}^{(2)}
=
1-(1-m_{E_1})(1-m_{E_2}).
\label{eq:depth_two_example}
\end{equation}

The boundary exploits consequently change from $\{E_1,E_2\}$ to $\{E_2,E_3,E_4,E_5\}$. Exploits $E_3$, $E_4$, and $E_5$ are first aggregated at $P_2$ and then propagated through $E_1$; they are not treated as direct causes of $P_1$. The number of non-zero boundary assignments correspondingly increases from $2^2-1=3$ to $2^4-1=15$.

The example generalizes recursively. An unexpanded boundary exploit contributes
\begin{equation}
m_E=\rho_Ev_E.
\label{eq:boundary_message}
\end{equation}
Under the MulVAL rule structure adopted in this work~\cite{ou2005mulval,ou2006scalable}, each exploit node has at most one privilege prerequisite. Accordingly, expanding an exploit $E$ that requires privilege $P$ replaces its boundary message with
\begin{equation}
m_E=\rho_Em_P,
\label{eq:expanded_message}
\end{equation}
provided that its non-privilege prerequisites are satisfied; otherwise, $m_E=0$. A privilege node combines its incoming exploit messages using
\begin{equation}
m_P
=
1-\prod_{E\in child(P)}(1-m_E).
\label{eq:privilege_message}
\end{equation}
Applying Eqs.~\eqref{eq:expanded_message} and \eqref{eq:privilege_message} recursively determines the target probability $q_P^{(d)}$ for any selected expansion depth.

\begin{proof}[Proof of Theorem~\ref{thm:depth_bp}]
Consider BP on the computation tree defined by the same dependency expansion, and let $b_z$ denote the BP message propagated from node $z$ toward the target privilege. Under the conditional-independence assumption used in Eq.~\eqref{eq:prob}, a boundary exploit sends
\begin{equation}
b_E=\rho_Ev_E.
\end{equation}
Thus, $b_E=m_E$ for every boundary exploit.

Assume that the BP messages and the recursively computed messages agree for the inputs of an internal node. For a privilege node, BP applies the noisy-OR update
\begin{equation}
b_P
=
1-\prod_{E\in child(P)}(1-b_E).
\end{equation}
Since $b_E=m_E$ for all incoming exploits, Eq.~\eqref{eq:privilege_message} gives $b_P=m_P$.

For an exploit node requiring privilege $P$, BP propagates
\begin{equation}
b_E=\rho_Eb_P
\end{equation}
when its non-privilege prerequisites are satisfied, and zero otherwise. Since $b_P=m_P$, Eq.~\eqref{eq:expanded_message} gives $b_E=m_E$.

Therefore, equality of the messages is preserved from the boundary to the target privilege. By structural induction over the expanded dependency tree, the resulting target probability $q_P^{(d)}$ equals the BP belief obtained from the same expansion.

For a cyclic dependency structure, the exploit-state vector consists of the exploit nodes within the cycle and the ordinary boundary exploit nodes on the remaining acyclic branches. For each non-zero assignment of this combined state vector, the target privilege probability is computed directly using loopy BP and used to construct the corresponding risk-clause weight.

When the expansion includes all messages used in the BP evaluation, the encoded privilege probabilities equal the BP beliefs used to evaluate $\mathcal{U}(\mathbf{x})$. Substituting these probabilities into the encoded risk terms therefore yields the same BP-based utility evaluation.
\end{proof}

Each expansion replaces a boundary exploit with its upstream exploit alternatives, increasing the number of boundary-state combinations. Since the WCNF encoding generates one risk clause for every non-zero boundary assignment, the number of risk clauses grows exponentially with the number of boundary exploits, which typically increases with the expansion depth. We therefore adopt the depth-one formulation in Eq.~\eqref{eq:prob} to balance the fidelity of the BP-based utility approximation and MaxSAT scalability.

% that's all folks
\end{document}

